\chapter{Protected Dehn disks and annuli}
\label{ch:dehn-surfaces}

The topological smoothing argument needs surfaces whose boundaries are
already marked in a particular handle.  Replacing a nullhomotopic circle
by some nearby embedded circle does not suffice: the final disk must take
every prescribed boundary value, and its interior must miss the whole old
boundary.  The annulus must retain both boundary maps.  Moreover, the
surfaces must enclose the same protected core that later changes of PL
structure are required to preserve.  We develop these requirements
separately, before combining them in the handle construction.

The classical starting points are the loop theorem and its Dehn-lemma
corollary, and the extension to planar surfaces by Shapiro and Whitehead.
Hamilton uses the protected cases on printed page 67 of
\cite{Hamilton1976}.  His short application does not contain all the
relative constructions needed here.  In particular, the marked tower,
the surgery of its projected surfaces, and the restoration of their
parameters are parts of the argument, rather than extra hypotheses.

\section{PL domains and full boundary markings}
\label{sec:dehn-domains}

Throughout this chapter put
\[
 E^k=\mathbb R^k,\qquad
 D^k=[-1,1]^k,\qquad Q^{k-1}=\partial D^k.
\]
Norms on these coordinate spaces are supremum norms.  Thus $D^2$ is a
square and $Q^1$ its polygonal boundary; these choices make literal
affine boundary formulas possible.  They are topological disks and
spheres, but no round Euclidean parametrization is being substituted.
Write $\partial R$ for the ambient frontier of a domain $R$, and
$\operatorname{int}_R$ for interior in the relative topology of $R$.

\begin{definition}[PL domain and finite PL maps]
\label{def:dehn-pl-domain}
Let $X$ be Hausdorff and let $e_i:U_i\to V_i\subset E^3$ be an open
atlas covering $X$, with locally piecewise affine transition maps in
both directions.  A PL domain is a closed subset $R\subset X$ such
that each $x\in\partial R$ has an atlas-compatible chart $b$ and an
affine functional $\lambda$ for which
\[
 \lambda(b(x))=0,\qquad
 y\in R\ \Longleftrightarrow\ \lambda(b(y))\geq0
 \quad(y\in\operatorname{dom}b).
\]
The linear part of $\lambda$ takes the value $1$ on some vector, so
this is a genuine halfspace chart.  No compactness or connectedness of
$R$ is included in this definition.

A finite PL map on a finite polyhedron is a map affine on the faces of
a finite subdivision.  For maps into $X$, this condition is imposed
locally in the supplied charts, with finite polyhedral witnesses on
the source.  A finite PL homeomorphism has finite PL forward and inverse
maps.  An embedded proper PL disk is an embedding $j:D^2\to R$ of this
kind with
\[
             j^{-1}(\partial R)=Q^1.
\]
Here properness includes this equality with the entire frontier, not
only avoidance of one chosen boundary patch.
\lean{PoincareMT.M76.PLDomain}
\source{Hamilton1976}
\dcref{Lemma 2 and its application, pp. 65--67}
\end{definition}

A finite triangulated surface used below has triangular maximal faces,
one incident triangle at each boundary edge and two at every other
edge, with interval or circle links at vertices.  Its Euler
characteristic is $v-e+f$.  These local incidence requirements are
essential when recognizing a disk or annulus from an Euler count.
The equation $v-e+f=0$ alone does not characterize an annulus.

Fix once and for all a once-around path $\omega$ on $Q^1$, with initial
point $q_*$.  If $F\subset\partial R$ is relatively open, a marked
singular disk consists of a map $f:D^2\to R$, a rim
$\gamma:Q^1\to F$ with $f|_{Q^1}=\gamma$, a base point $z\in F$,
and a path $p$ in $F$ from $z$ to $\gamma(q_*)$.  For a normal subgroup
$J\triangleleft\pi_1(F,z)$ its exclusion condition is
\[
            [p*(\gamma\circ\omega)*p^{-1}]\notin J.
\]
A homotopy of rims changes $p$ by appending the trace of $q_*$.
This convention keeps every exclusion statement in the original group.
The marked loop theorem may change $\gamma$; the proper Dehn theorem
will subsequently restore the original $\gamma$ pointwise.

\section{Finite towers and marked disk descent}
\label{sec:dehn-marked-tower}

\begin{theorem}[Marked boundary loop disk]
\label{thm:dehn-marked-loop}
\uses{def:dehn-pl-domain}
Let $R$ be a PL domain in a Hausdorff space with the atlas of
Definition~\ref{def:dehn-pl-domain}.  Let $F\subset\partial R$ be
relatively open, $f:D^2\to R$ continuous, and $\gamma:Q^1\to F$
continuous with $f|_{Q^1}=\gamma$.  For every normal subgroup
$J\triangleleft\pi_1(F,\gamma(q_*))$ excluding $[\gamma\circ\omega]$,
there are a proper embedded PL disk $j:D^2\to R$, a continuous rim
$\gamma':Q^1\to F$, and a path $p$ from $\gamma(q_*)$ to
$\gamma'(q_*)$ in $F$ such that
\[
 j|_{Q^1}=\gamma',\qquad
 [p*(\gamma'\circ\omega)*p^{-1}]\notin J.
\]
\lean{PoincareMT.M76.Dehn.exists_marked_boundary_PL_loop_disk}
\source{Hatcher3M2014}
\dcref{Section 3.1, Theorem 3.1 and the double-arc calculation, pp. 54--57}
\end{theorem}

We describe the tower and the two separate termination arguments.  A
finite PL approximation of $f$ is made from the compact graph model
around its image.  Finitely many chart coordinates give a continuous
graph map $G$ which is locally piecewise affine in every original chart;
a compact set $C$ contains the image in its interior.  The scalar
defining function for the original half-space is retained as the last
graph coordinate.  A finite polyhedral pair $(P,Z)$ is then chosen with
\[
 P=G(C\cap R),\qquad Z=G(C\cap\partial R),
\]
and the literal inverse $H:(C\cap R)\simeq P$ transports a finite PL
approximation in $P$ back to $R$.  The approximation theorem supplies a
homotopy in $R$ whose tracks over the marked source subset stay in $F$;
it is relative only to that carrier, so the rim parametrization may
change.  The inverse identities for $G$ and $H$ preserve the whole
frontier equation, rather than only a chosen boundary patch, and the
trace at $q_*$ transports the excluded class to the new rim.

A tower stage is a connected neighborhood of the lifted singular
surface, with its induced PL atlas, a local-homeomorphism projection
to the original neighborhood, and a deformation retraction onto the
lifted image.  A defining function $r$ records the original domain
and frontier as $r\geq0$ and $r=0$.  The deformation preserves these
sets.  A successor stage is constructed from an actual connected
two-sheeted cover of the preceding stage's carrier by replacing it
with the connected model consisting of that carrier times
$\{0,1\}$ and its transported covering topology, rather than the product topology.
The source lift is PL in the lifted chart cover.  Its image is path
connected, and the PL deformation theorem supplies the open level set
\[
 O=\{x:z(x)>1/2\}
\]
with a deformation retraction onto that image.  Restricting the lifted
charts to $O$ gives the successor; projection and inverse-chart
identities are retained on their complete domains.  The defining
function is nonnegative (and zero) along the deformation exactly when
it was so at the preceding stage, so the full old domain and frontier
remain available.  For a disk the source lift exists because the source
is simply connected; the annulus tower uses the same construction only
after choosing a closed lift of its marked negative rim.

The ascent has a finite measure independent of newly chosen
triangulations.  Here is the construction of the fixed source marks.
Write the initial graph-coordinate map as $F$, affine on each of
finitely many compact convex cells $C_i$, with affine expression $A_i$.
For every nonempty intersection $A_i(C_i)\cap A_j(C_j)$ choose one
point and a preimage in each cell.  Let $T$ consist of all these
chosen source points.  This choice is made before any cover is chosen.

Suppose a later lift $g$ has a locally injective graph projection $q$
with $qg=F$.  On a fiber of $A_i|_{C_i}$ the map $g$ is constant:
that fiber is convex and connected, whereas local injectivity of $q$
makes its fiber discrete.  The compact-to-Hausdorff quotient
$C_i\to A_i(C_i)$ therefore factors $g$ continuously through a map
$g_i:A_i(C_i)\to\operatorname{target}(g)$.  Two such local sections
of $q$ agreeing at one point agree on the whole convex overlap.
Their equality set is closed by Hausdorffness and open by local
injectivity, so connectedness proves the assertion.  The same
argument applies to $g_i$ and a deck translate of $g_j$.

In a connected double cover let $\tau$ be its fixed-point-free deck
involution.  The inverse image of the preceding source image is
$g(D^2)\cup\tau g(D^2)$.  Lifting the strong deformation retraction
shows that this inverse image is connected.  Its two nonempty
compact, hence closed, pieces must meet.  An intersection lies over
some pair of original cells; the preceding overlap argument moves
it to their chosen marks $a,b\in T$.  Thus
$g(a)=\tau g(b)\ne g(b)$, whereas their projections at the preceding
stage agree.  Projection maps the new finite marked image onto the
old one and identifies at least this pair.  Its cardinality therefore
strictly decreases under projection.

If $s$ is a stage and $s_T$ its source map on $T$, use
\[
                    m(s)=|T|-|s_T(T)|.
\]
The preceding argument proves that $m$ strictly decreases at every
successor, including restriction to the next neighborhood, whose
inclusion in the cover is injective.  Taking a stage
of minimal measure among all stages reachable by finite towers gives a terminal
stage, without assuming a choice of covers in advance.

For disks terminality excludes every connected double cover.  The
compact relative neighborhood has a finite PL model with its actual
frontier as a subcomplex.  A nonzero mod-two edge cocycle modulo vertex
coboundaries would give a connected two-sheeted cover by the rule
\[
       (v,\epsilon)\longmapsto(w,\epsilon+c(vw)).
\]
The triangle cocycle equation makes the lifts agree on triangles;
changing a vertex potential changes the sheet labels.  Terminality
therefore forces first mod-two cohomology to vanish.  The boundary
calculation below then makes first mod-two cohomology of the whole
frontier vanish.  Each connected closed frontier component has
Euler characteristic two and the circle links of a triangulated
surface; finite PL sphere recognition applies.
In a finite sphere model the marked loop is decomposed into simple
polygonal loops.  If all their based classes belonged to $J$, the
original product would belong to $J$; select one that does not.
It bounds a finite PL disk on the sphere.  Push its interior inward in
a collar, keeping the rim on the original marked part of the frontier.
This constructs a terminal proper marked disk.

Descent is not accomplished by merely projecting an embedding: the
projection may identify its two sheets.  After finite relative
normalization, each collision is an ordinary transverse double curve.
In local coordinates its branches are the two planes
$y=x$ and $y=-x$, with the appropriate halfspace at an endpoint on
the frontier.  There are no triple values because the covering has
two sheets and the upper disk is embedded.  Compactness, finite PL
intersection data, and local injectivity yield a finite double locus:
interval components ending on the rim and circle components in the
interior.  Normalization is implemented by a finite sequence of
supported motions.  In each disjoint exceptional window the complete
composite has exactly the branch images of the corresponding single
repair; outside their closed union it has the original branch images.
Consequently the local crossing assertion is about whole branch
images, not just their tangent planes or germs.

The finiteness here is a scheduled composite, not an appeal to generic
position.  The exceptional source pairs are covered by a finite list of
compact repair windows with pairwise disjoint projected supports.  The
composite is the identity off the closed union of those supports; inside
one window its projected branch image agrees with the corresponding
single repair, so the supplied crossing chart transports through the
composite.  The exceptional pairs are therefore all repaired by the
finite list, while pairs outside the union retain the original crossing
chart.  The same supported homeomorphisms preserve the domain and the
marked frontier, and their time-dependent composite gives the rim
homotopy used when descending one tower step.

For an interval component choose a closed transverse prism whose
inverse image consists of its two source strips.  Shrink its width
uniformly so that its frontier ends lie in $F$.  The prism is built
along the finite polygonal axis: interior blocks meet in prescribed
faces, and the endpoint half-blocks retain the old frontier.  Cutting
the two strips and joining the alternative sides gives two candidate
disks.  Their finite PL maps agree on the full joining intervals.
More explicitly, normalize the closed prism to
$[-1,1]^2\times[0,1]$, with coordinates $(x,y,s)$ and original sheets
$y=x$ and $y=-x$.  For $0<b<1$ put $H_b(u)=\max(|u|,b)$.
The two resolutions have transverse arcs
\[
 (u,H_b(u)),\ (u,-H_b(u)),\qquad\hbox{or}\qquad
 (H_b(u),u),\ (-H_b(u),u).
\]
Keep $s$ unchanged.  These are finite PL strips, affine after
subdivision at $u=\pm b$, and each has a coordinate projection as
left inverse.  The positive and negative strips within a pairing
are disjoint, separated by at least $2b$.  Where $|u|\geq b$
their images are the old crossing arms, so the entire outer collar
and all four corner lines have their original values.  Preservation
of $s$ also preserves both frontier end planes.

Removing the interiors of the two source strips from the old disk
leaves three finite PL disks, denoted $A,M,C$ in their planar order.
The signs of the middle-facing arms specify a reflection or
interchange of the transverse prism coordinates.  In those coordinates
the first candidate is the chain consisting of $A$, one upper
replacement strip, and $C$; the second is the chain consisting of
$A$, a left replacement strip, $M$, a right replacement strip,
and $C$.  Each attachment is along a complete boundary interval
with the same parametrization on both sides.  Triangulate these
intervals simultaneously and extend their subdivisions over the
adjacent disks.  Pasting the affine formulas constructs each map
on a finite PL disk.  This construction retains the old map on
every retained disk, and gives an exact formula for the inverse
image of every target set as the union of its inverse images in
these pieces.  Applying it to the tube and to the old double locus
is what justifies the collision and count assertions below.

The choice between the candidates is algebraic.  Write $a,b,c,d$ for
the four successive old rim paths, transported to one base point by
the paths in the endpoint squares.  In the first pairing the candidate
words are $u=ac$ and $v=ab^{-1}cd^{-1}$, and
\[
             abcd=u\,d^{-1}v^{-1}u\,d.
\]
In the second pairing they are $u=ac^{-1}$ and $v=adcb$, and
\[
             abcd=u(cd)^{-1}u^{-1}v(cd).
\]
Both identities follow by cancellation in the fundamental group.
Thus if both candidates lay in the normal subgroup $J$, so would the
old word.  Select a candidate outside $J$.  Actual source boundary
traversals can be inverses or cyclic rotations of these words; inverse
closure and normality justify precisely these changes, with their
basepoint paths retained.

The interval surgery strictly reduces the number of source double
components meeting $Q^1$, and does not increase the number of interior
components.  Strong induction removes every boundary component.  The
rim is then embedded: two distinct rim points with the same value
would lie in the now absent boundary double locus.  Interior circle
surgery subsequently fixes the whole rim and strictly reduces the
interior count.  Paired circles are resolved using their two source
annuli; a self-paired circle uses a reflected tube.  The source disks
and annuli retained in each resolution identify its remaining double
components with old components away from the selected one.  At zero
count the whole disk is injective.  Compactness of $D^2$ and the
Hausdorff property of the target turn this injection into an embedding.
Backward induction through the finite tower, followed by its initial
open embedding, proves Theorem~\ref{thm:dehn-marked-loop}.

\begin{lemma}[Boundary rank from finite incidence]
\label{lem:dehn-boundary-rank}
\uses{def:dehn-pl-domain}
Let $K$ be a finite triangulation of a connected compact PL
three-dimensional domain with nonempty boundary subcomplex $A$.
Use mod-two chains and write $b_i$ for their homology dimensions.
Then
\[
                         b_1(A)\leq2b_1(K).
\]
The assertion allows disconnected $A$.
\lean{AbstractSimplicialComplex.boundary_incidence_rank_le_of_counts}
\end{lemma}

\begin{proof}
Put $v,e,f,t$ for the numbers of vertices, edges, triangles and
tetrahedra of $K$, and $v_A,e_A,f_A$ for the boundary counts.
An interior triangle has two incident tetrahedra and a boundary
triangle has one.  The link of an interior vertex is a sphere of
Euler count two; the link of a boundary vertex is a disk of count
one.  Summing these vertex-link counts gives
\[
 2e-3f+4t=2v-v_A,
 \qquad 4t=2f-f_A,
 \qquad 3f_A=2e_A.
\]
Elimination yields the literal incidence identity
\[
 2v+2f+e_A=2e+2t+v_A+f_A,
 \quad\hbox{or equivalently}\quad \chi(A)=2\chi(K).
\]
The link models here come from the original interior and halfspace
charts after subdivision; the boundary is the entire frontier,
not a chosen subset of faces.

Let $c$ be the number of components of $A$.  A boundary two-cycle
has the same coefficient on the triangles adjacent to any edge.
Propagation across triangle adjacency shows that it is constant
on each component.  Thus its cycle space has dimension $c$.
The tetrahedron adjacency graph of $K$ is connected.  Consequently
a sum of tetrahedra whose boundary is supported in $A$ has the
same coefficient on every tetrahedron: interior triangular faces
force equality across adjacent tetrahedra.  Its boundary is either
zero or the sum of all boundary triangles.  Since $A$ is nonempty,
the tetrahedron boundary map is injective and its image has dimension
$t$.  Its intersection with the $c$-dimensional space of boundary
two-cycles has dimension exactly one.  Both spaces lie in $Z_2(K)$,
so
\[
                       c+t-1\leq\dim Z_2(K).
\]
Connectedness gives $\operatorname{rank}\partial_1=v-1$.
Writing $b=b_1(K)$, rank-nullity for $\partial_2$ now gives
\[
 \operatorname{rank}\partial_2=e-v+1-b,
 \qquad \dim Z_2(K)=f-e+v-1+b=\chi(K)+t-1+b.
\]
Hence $c\leq\chi(K)+b$.  On the closed surface $A$, its zeroth
and second Betti numbers both equal $c$, so
\[
 b_1(A)=2c-\chi(A)=2c-2\chi(K)\leq2b.
\]
This proves the bound directly from the finite incidence maps.
It is the same calculation for $b=0$ in the disk tower and $b\leq1$
in the annulus tower.
\end{proof}

\paragraph{Whole-image control.}
\label{rem:dehn-local-review}
The ordinary crossing prisms, their source strips, and the terminal
chart stars supply geometric hypotheses in the preceding proofs.
Their whole-image and whole-frontier properties matter: a collection
of local pictures without those equalities would not justify descent.
The following sections give the
replacement maps and the boundary correction separately.

\section{Restoring a prescribed disk boundary}
\label{sec:dehn-proper-disk}

\begin{theorem}[Proper disk with the original parametrization]
\label{thm:dehn-standard-proper}
\uses{thm:dehn-marked-loop}
Let $R\subset E^3$ be compact and a PL domain for the standard affine
atlas.  Let $S\subset\partial R$, let $\gamma:Q^1\to S$ be a finite
PL homeomorphism, and let $f:D^2\to R$ be continuous with
$f|_{Q^1}=\gamma$.  There exist a compact $D\subset R$ and a finite
PL homeomorphism $b:D^2\to D$ such that
\[
 b|_{Q^1}=\gamma,\qquad
 b(x)\in\partial R\ \Longleftrightarrow\ x\in Q^1.
\]
\lean{PoincareMT.M76.Dehn.exists_standard_proper_disk}
\source{Hatcher3M2014}
\dcref{Section 3.1, Corollary 3.2, printed p. 57}
\cite[Corollary 3.2, p. 57]{Hatcher3M2014}
\end{theorem}

\begin{proof}
Construct a finite PL annular neighborhood $T$ of $S$ in $\partial R$,
with $S$ its middle circle, and a relatively open mark $F$ containing
$S$ and lying strictly between the outer rims of $T$.  Local boundary
halfspace charts and a subdivision of the embedded polygonal circle
give this neighborhood.  The inverse middle-circle parametrization
followed by radial retraction shows that the once-around loop of
$\gamma$ is nontrivial in $T$.

Apply Theorem~\ref{thm:dehn-marked-loop} with
\[
 J=\ker\bigl(\pi_1(F,\gamma(q_*))\longrightarrow
                    \pi_1(T,\gamma(q_*))\bigr).
\]
The resulting embedded proper disk has a rim essential in the whole
annulus $T$.  Using this kernel is necessary: nontriviality in a
smaller open mark alone would not give essentiality in $T$.

Pull the output rim back through the finite PL annulus chart.
Subdividing its parametrization gives a simple polygon.  Its bounded
planar side contains the inner square.  Indeed, otherwise a point of
the inner square would lie in the unbounded complementary component,
and the polygon would have winding zero in the annulus, contradicting
essentiality.  Its closure lies strictly inside the outer square;
the polygon lies in that convex square and its bounded side cannot
cross the square's supporting lines.  Polygonal annulus recognition
now gives two finite PL annuli $A_0,A_1$ in $T$ with the same outer
rim $Q$: their inner rims are respectively the output rim $S_0$ and
the desired rim $S_1=S$.  They may overlap.

Choose a finite PL inward collar $c(z,t)$ of the whole frontier, with
$0\leq t\leq1$ and $c(z,0)=z$.  We construct the compression on the
whole finite polyhedral model of $R$, including the part outside the
collar.  Write
\[
 C_c=c(\partial R\times[0,1]),\qquad
 O_c=c(\partial R\times[0,1)),\qquad
 B_c=c(\partial R\times\{1\}).
\]
Here $O_c$ is relatively open in $R$.  The closed strip $C_c$ and
its roof $B_c$ are finite polyhedra, and collar injectivity gives
$C_c\setminus O_c=B_c$.  Subdivide the finite models of $R$ and
$C_c$ together so that the latter is a subcomplex.  The closures
of the remaining simplex interiors form a finite polyhedron with
carrier $\overline{R\setminus C_c}$.  Adjoining the roof and taking
a common subdivision therefore gives a finite residual $P_c$ with
\begin{equation}
 \begin{aligned}
 P_c&=\overline{R\setminus C_c}\cup B_c=R\setminus O_c,\\
 C_c\cap P_c&=B_c,\qquad C_c\cup P_c=R.
 \end{aligned}
 \label{eq:dehn-collar-residual}
\end{equation}
For the middle equality, $R\setminus O_c$ is closed and contains
$R\setminus C_c$ and $B_c$; conversely each of its points either
lies outside $C_c$ or belongs to $C_c\setminus O_c=B_c$.

On $C_c$ use $c(z,t)\mapsto c(z,(t+1)/2)$, and on $P_c$ use
the identity.  These finite PL maps agree on their entire overlap
$B_c$.  The shifted strip meets $P_c$ exactly on that same roof,
so the pasted map is injective.  Compactness makes it a homeomorphism
onto its image $R'\subset R$; the inverses on its finitely many
affine pieces show that it is a finite PL homeomorphism.  Its image satisfies
\[
 c(z,t)\in R'\ \Longleftrightarrow\ t\geq\tfrac12.
\]
The compression carries the old disk to a disk $B_0\subset R'$ and
its rim to $c(S_0,1/2)$.  Here the collar is constructed on a finite
boundary model; $z$ denotes its boundary coordinates.
Parametrize each $A_i$ by a square annulus of depth $d\in[-1,1]$,
where $d=1$ is its inner rim.  Lift the two annuli into the collar
using the heights
\[
                 h_0(d)=\frac{3+d}{8},\qquad
                 h_1(d)=\frac{1-d}{8}.
\]
These maps are finite PL and injective, since their base annulus
maps are injective and the collar is injective.  Their ranges of
heights are $[1/4,1/2]$ and $[0,1/4]$.  Equality between points on
the two graphs forces both depths to be $-1$ and their base values
to be the same point of $Q$.  Thus their intersection is exactly
$c(Q,1/4)$.  The first graph meets $R'$ exactly on
$c(S_0,1/2)$; the second misses $R'$.  Gluing $B_0$ to the first
graph and then to the second therefore produces an embedded finite
PL disk, regardless of overlaps of $A_0,A_1$ on the original frontier.
Only the second graph has height zero, precisely on $S_1$.
Hence the new disk meets the entire old frontier exactly in $S$.

Finally choose any finite PL disk parametrization of this image.
Its boundary map differs from $\gamma$ by a finite PL homeomorphism
of $Q^1$.  Subdivide the circle so that this map is affine on each
edge, triangulate the disk by coning those edges to an interior
vertex, and extend linearly over each cone.  The inverse is obtained
from the inverse boundary subdivision.  Composing this extension
with the chosen disk parametrization restores every value of
$\gamma$ and preserves properness.
\end{proof}

\begin{lemma}[Confinement to one chart]
\label{lem:dehn-chart-proper}
\uses{thm:dehn-standard-proper}
Let $R$ be a PL domain in a Hausdorff $X$, let $e_i$ be one of its
charts, and let $U\ne\varnothing$ be an open subset of its target.
Suppose a finite PL embedded coordinate rim $\gamma:Q^1\to U$
has a continuous filling in
$\{y\in U:e_i^{-1}(y)\in R\}$ and
$e_i^{-1}(\gamma(Q^1))\subset\partial R$.
There is a compact proper chartwise PL disk contained in
$R\cap e_i^{-1}(U)$, with boundary $e_i^{-1}\circ\gamma$ pointwise.
\lean{PoincareMT.M76.Dehn.exists_proper_disk_in_chart}
\end{lemma}

\begin{proof}
The filling image $A$ is compact.  We first give the relative compact
neighborhood construction in a PL domain $R_0$ in an open ambient
space.  Given an open $W\supset A$, choose finitely many relatively
compact chart neighborhoods covering $A$, with closures in $W$.
In each chart a finite PL cutoff is one on a smaller neighborhood
and zero outside a compact subset of the larger one.  Extending by
zero and taking the finite maximum gives a continuous chartwise PL
function $w:X\to[0,1]$, equal to one on an open neighborhood $V_0$
of $A$, with compact support in $W$.

The band $1/3\leq w\leq2/3$ is compact.  Cover its old-boundary
part by finitely many halfspace charts and triangulate them so that
$w$ is affine on faces and the old boundary is a subcomplex.  Avoid
the finitely many vertex heights in these models.  Inside a remaining
open interval choose an ambient regular PL level $t$ by the same
finite vertex-height avoidance.  Thus $1/3<t<2/3$, and
\[
                        K=R_0\cap\{w\geq t\}
\]
is compact and contained in $R_0\cap W$.  At an old-boundary point
where $w>t$, its frontier has the old halfspace chart.  At a level
point in the interior of $R_0$, the height chart gives a halfspace.
At an intersection of the two faces, the relative PL corner chart
straightens their quadrant to a halfspace while distinguishing the
old face and the new face on its boundary.  The nonvertex choice on
the boundary triangulation is required for this last chart; ambient
regularity alone would not suffice.  These three cases give PL-domain
charts along the whole frontier of $K$.

On $V=V_0\cap W$ one has $K\cap V=R_0\cap V$.  As $V$ is open,
frontiers also agree there.  In particular $A$ is retained, including
all its old-boundary points, and lies in the relative interior of
$K$ in $R_0$.  For the lemma apply this construction to the coordinate
domain in $U$ and to the compact filling image.  The compact set $K$
is closed in $E^3$ as well as in $U$; the open inclusion transports
its boundary charts to standard ambient charts.  Apply
Theorem~\ref{thm:dehn-standard-proper} to $K$.
If an interior point of the resulting disk lay on the original
frontier, it would also lie on the frontier of the compact subdomain:
an interior point of the latter is an interior point of the former
by containment.  This contradicts disk properness.  The chart and
the open inclusion preserve frontier membership locally, so the
inverse chart returns the asserted disk and its full boundary map.
The boundary of $U$ need not itself be PL.
\end{proof}

\section{The protected handle and its two disk slabs}
\label{sec:dehn-protected-disks}

Fix $(k,n)=(1,2)$ or $(2,1)$ and a discrete subgroup
$L\subset E^n$.  Put
\[
 X=E^k\times(E^n/L),\quad R=D^k\times(E^n/L),\quad
 \pi(x,y)=(x,[y]),
\]
\[
 C_a=\pi(D^k\times aD^n),\qquad
 B_a=\pi(Q^{k-1}\times aD^n).
\]
The protected compact set is $A=C_1$.  Assume $R$ has a supplied
ambient PL atlas $e$, and let $h$ be a topological open coordinate
map into $E^3$ whose domain contains $D^k\times E^n$.  Assume an
open neighborhood $N$ of $\partial(D^k\times E^n)$ is given and
$h$ is locally PL on $N\cap\operatorname{dom}h$.
There is no PL assumption on the interior of $h$.

\begin{definition}[Retained block]
\label{def:dehn-retained-block}
\uses{def:dehn-pl-domain}
A retained block means that $\pi$ is injective on $D^k\times2D^n$
and that one specified chart $e_*$ contains $C_2$ and satisfies
\[
                 e_*(\pi(x,y))=h(x,y)
                 \quad ((x,y)\in D^k\times2D^n).
\]
All subsequent surfaces use this same lattice, projection and chart.
\lean{PoincareMT.M76.HamiltonRetainedBlockChart}
\source{Hamilton1976}
\dcref{Construction for indices one and two, printed p. 67}
\end{definition}

\begin{theorem}[Simultaneously disjoint protected disks]
\label{thm:dehn-protected-disks}
\uses{def:dehn-retained-block,lem:dehn-chart-proper}
For $k=2,n=1$, the preceding data produce two compact embedded
chartwise PL disks $j_\pm:D^2\to R$ with disjoint images $T_\pm$,
satisfying
\[
 j_\pm(u)=\pi(u,\pm3/2)\quad(u\in Q^1),\qquad
 j_\pm^{-1}(\partial R)=Q^1,
\]
\[
 T_-\subset\pi(D^2\times(-2,-1)),\qquad
 T_+\subset\pi(D^2\times(1,2)).
\]
In particular each image lies in $C_2\setminus C_1$ and in
$\pi(D^2\times(-2,2))$.
\lean{PoincareMT.M76.Dehn.ProtectedDisks.exists_protected_disk_pair}
\source{Hamilton1976}
\dcref{Index-two generalized Dehn application, printed p. 67}
\end{theorem}

\begin{proof}
The maps $u\mapsto h(u,\pm3/2)$ are continuous disk fillings.  Their
rims are finite PL, because they lie in the compact old boundary
where $h$ is locally PL.  Use the open ambient slabs in $X$ with
transverse coordinate in $(-2,-1)$ and $(1,2)$, intersect them with
the retained chart, and apply Lemma~\ref{lem:dehn-chart-proper}.
The inverse chart gives exactly $\pi(u,\pm3/2)$ at every rim point.

The slabs are disjoint in $X$, not merely in source coordinates.
If their projected points coincided, both source points would lie
in $D^2\times2D^1$, and retained injectivity would identify their
transverse coordinates, which have opposite signs.  The same
injectivity proves avoidance of $C_1$: a point in either strict slab
cannot also have a representative with transverse norm at most one.
Thus independent disk constructions in these two specified open
slabs give the simultaneous disjointness required.
\end{proof}

\section{A marked annulus at the terminal stage}
\label{sec:dehn-terminal-annulus}

\begin{theorem}[Embedded annulus in the terminal tower stage]
\label{thm:dehn-terminal-annulus}
\uses{def:dehn-retained-block,lem:dehn-boundary-rank}
\lean{PoincareMT.M76.Dehn.ProtectedAnnulus.ProtectedAnnulusTerminalData.exists_embedded_terminal_annulus}
The constructed terminal stage of a protected annulus tower contains
a finite PL embedded annulus on the unchanged source complex. Its
image lies in the retained chart domain, its intersection with the
whole domain frontier is precisely the two source rims, and it agrees
pointwise with both prescribed lifted rim maps.
\end{theorem}

For $k=1,n=2$ the original continuous annulus in chart coordinates is
\[
       (t,u)\longmapsto h(t,\tfrac32u),
       \qquad (t,u)\in D^1\times Q^1.
\]
It lies in the strict shell $h(D^1\times\{1<\|y\|<2\})$.
On this shell the map
\[
             z\longmapsto
             \frac{(h^{-1}z)_2}{\|(h^{-1}z)_2\|}
\]
is a continuous left inverse on both prescribed rims.
Their once-around classes are consequently nontrivial.  Translation
in the $t$ coordinate gives an actual homotopy between the two rims.
These are properties of the specified continuous chart; no annulus
in the new PL structure has yet been constructed.

Approximate the cylinder by a finite PL cylinder relative to both
complete rims, preserving the equivalence between belonging to the
old frontier and being an endpoint $t=\pm1$.  Construct a tower as
in Section~\ref{sec:dehn-marked-tower}, but admit a double cover only
when the whole negative rim has a closed lift.  The cylinder lifts
by homotopy lifting from that rim.  The same finite-source measure
strictly decreases.  The terminal condition is therefore
\[
 \text{no connected two-sheeted cover admits a closed lift
 of the prescribed negative rim.}
\]
It is not the assertion that the stage has no double covers.

Take a compact relative PL neighborhood $K$ of the terminal singular
cylinder, with deformation retraction to its image and with a finite
triangulation of its whole boundary.  Refine it to include both marked
rims.  A mod-two cocycle whose restriction to the negative rim is
exact constructs a double cover in which that rim lifts closed.
Terminality forces that cocycle to be exact.  Restriction in first
cohomology is thus injective into the one-dimensional cohomology of
the circle, giving
\[
       \dim H^1(K;\mathbb F_2)\leq1,
       \qquad \dim H^1(\partial K;\mathbb F_2)\leq2.
\]
The second inequality is Lemma~\ref{lem:dehn-boundary-rank} for the
same compact region and its full boundary subcomplex.  The construction
retains compatible face orientations from the actual charts; this
also orients the closed boundary components.

Each boundary component containing a marked rim has a radial
retraction onto that rim.  Such a component has nonzero first
cohomology.  For a connected closed orientable surface the dimension
of first mod-two cohomology is even, so it is at least two.  If the
two rims were in different components, the total boundary dimension
would be at least four, contrary to the displayed bound.  They
therefore lie on one component, of first cohomology dimension two
and Euler characteristic zero.

Remove disjoint thin open annular neighborhoods of the two rims in
that component.  The closed cut surface has four boundary circles,
with all edge incidences inherited from a common subdivision.  Cutting
off annuli changes neither total Euler characteristic nor its value
zero.  None of these four offset rims bounds a disk in the cut surface:
such a disk, followed by the radial map, would contract a once-around
circle.  The finite surface argument now selects a connected
count-zero piece joining an offset rim from each original collar.
This is done on the literal cut carrier, not on an abstract surface
with the same Euler number: every edge has one triangular coface
precisely on one of the four collar circles and two elsewhere, and
every vertex link is connected.  These coface and link conditions make
the selected piece a finite PL surface with its displayed boundary
circles.  To see the topological reason, cutting a torus along an
essential circle gives an annulus; a second disjoint essential circle
lies in that annulus and separates its two boundary circles.  The two
complementary pieces are annuli, and either connects the two collars.
In the finite model cap the selected two boundary circles by cones.
The coface count becomes two at every interior edge, the capped links
are circles, and finite sphere recognition followed by removal of the
two open cone disks gives an annulus on the original carrier.  Thus
the output keeps both complete rim subcomplexes as sets before the
final collar correction.

Attach back the appropriate half of each collar.  This produces a
finite PL boundary annulus whose boundary sets are the original
marked rims.  It remains to correct the parametrizations.  After
identifying the rim circle with $\mathbb R/32\mathbb Z$, the two
required circle homeomorphisms have the same degree, since they are
homotopic.  Cut each at one point.  Its lift has the form
\[
       s\longmapsto a_\pm+\varepsilon32 e_\pm(s),
       \qquad 0\leq s\leq1,
\]
where $e_\pm$ are increasing finite PL interval homeomorphisms
fixing $0,1$, and the common sign $\varepsilon$ is $1$ or $-1$.
Triangulate the rectangle to extend the two interval boundary maps;
its compatible side identifications give an annulus homeomorphism.
A common reflection realizes $\varepsilon=-1$ if needed, and a PL
phase shear realizes the two independent shifts $a_\pm$.
Composition restores both whole circle maps.  A collar push of the
annulus interior, fixing its rims, then gives a proper embedded
annulus in the terminal region.

\section{Circle surgeries and annulus descent}
\label{sec:dehn-annulus-descent}

Project the terminal annulus down one tower step.  Finite face motions
and exceptional repairs can be supported off the entire marked rim.
They preserve the upper embedding and the domain.  Their projected
double locus is a compact finite polyhedron with ordinary crossing
charts, no triple points, and positive separation from the rim.
It is therefore a finite union of source circles.  Each double point
has one partner, and partnership permutes those circles by an
involution.  Denote the number of connected components of this
actual source double locus by $c(f)$.

We spell out how the ordinary local models yield a tube with the
full preimage required by surgery.  In a crossing chart the axis
is the intersection of two coordinate planes.  Refine the finite
graph-coordinate model so that the axis, the two planes and their
four signed sectors are subcomplexes, and so that the closed star
of each axis vertex lies in one crossing chart.  Barycentric dual
blocks of consecutive axis edges give transverse disks; blocks of
axis vertices give intervening prisms.  Parametrize each transverse
disk by a diamond divided by its two axes into four triangles.
The four arms inherit the two sheet labels and their side signs.

On a vertex block, first parametrize its axis interval and its four
sheet half-faces, matching the transverse disks at both ends.
Each closed sector is a PL ball with that prescribed boundary
decomposition.  Extend the boundary parametrization over the sector
by a cone subdivision.  Adjacent sectors have exactly one prescribed
half-face in common; opposite sectors meet exactly on the axis.
Their maps consequently paste injectively into a map of the entire
diamond prism.  This proves full sector, sheet, axis and end-face
inverse-image equalities.  Successive blocks meet on the same
transverse disk maps.  Composing their signed arm permutations
therefore constructs a tube, together with its closing permutation
when the axis is a circle.  Compactness permits a common positive
shrinking of the diamonds within these charts.  The compact source
outside the two local branch neighborhoods has image avoiding the
selected axis; shrink once more to exclude that image.  The whole
tube preimage is now precisely the source strips obtained by
inverting the two sheet maps, with no third source pieces.

For a circular axis the closing could a priori permute and reverse
the two axes.  Planarity of the source excludes a reversal of a
separately closed strip: a point on one side of its middle polygon
would return on the other side without meeting the polygon.
Thus two separately closed source circles give the identity closing.
When the strips exchange, the same inside-outside argument forces
their two signs to agree; a quarter turn or a single sign reversal
is impossible.  Changing the initial signed frame converts this
remaining closing to reflection of one transverse coordinate.
These are precisely the identity and reflected tubes used below.
For an interval axis the same block construction has halfspace
end blocks, retaining the complete old-frontier end faces; there is
no circular closing condition.

\begin{lemma}[Decreasing annulus surgery]
\label{lem:dehn-annulus-decrease}
\uses{def:dehn-retained-block}
For a proper finite PL projected annulus in one of these retained
stages, with the prescribed two rim maps, ordinary crossing charts
and nonempty double locus, there is another such annulus map $g$
with the same whole rim maps and
\[
                           c(g)<c(f).
\]
The output includes a finite ordinary circle decomposition, and all
its double values are interior points of the same domain.
\lean{PoincareMT.M76.Dehn.ProtectedAnnulus.exists_decreasing_planar_stage_annulus}
\end{lemma}

The construction has three cases.  A circle equal to its partner has
a reflected transverse tube.  Write it as the image of
$[-d,d]^2\times[0,2L]$, where $d>0$ and $4d<L$, under a map $\tau$
whose only identifications are
\[
                   \tau(x,y,0)=\tau(x,-y,2L).
\]
The old source collar has period $4L$: its first half has value
$\tau(u,u,s)$ and its second half has value
$\tau(u,-u,s-2L)$.  For $0<b<d$ an embedded resolving annulus is
given, before boundary correction, by
\[
 a(s,u)=
 \begin{cases}
 \tau(u,H_b(u),s),&0\leq s\leq2L,\\
 \tau(u,-H_b(u),s-2L),&2L\leq s\leq4L.
 \end{cases}
\]
The fiber law pastes the two period seams.  Away from those seams,
the first and second halves are separated by the nonzero sign of
the second coordinate.  Within a half, that same law recovers
$u$ and $s$.  Thus the only identifications are those already
made in the source annulus.  The graph meets the side of the tube
exactly at $|u|=d$, since $b<d$.  Subdivision at $u=\pm b$ gives
finite PL formulas.

At $u=d$ these are the old boundary values.  At $u=-d$ they differ
by half a period.  Precompose by the annular shear
\[
                 (s,u)\longmapsto
                 \bigl(s+L-(L/d)u\pmod {4L},u\bigr).
\]
Its shift is zero on $u=d$ and $2L$ on $u=-d$; its inverse subtracts
the same function of $u$.  The square-annulus strip coordinates
realize this as a finite PL homeomorphism after subdivision.  The
corrected resolution consequently agrees with the original map at
every point of both collar boundary circles.

Let $T$ be this source collar and $C$ the old source minus its
interior.  Paste the corrected resolution on $T$ to the unchanged
map on $C$.  If a new point of $T$ and an old point of $C$ have the
same image, the exact old tube preimage puts the latter point in
$C\cap T=\partial T$.  Boundary agreement and injectivity of the
resolution then identify the source points.  Hence every point of
$T$ has a unique new image preimage, and the entire new double
relation is the old double relation restricted to $C\times C$.
The selected middle circle has been removed.  All remaining
crossings have open neighborhoods in the unchanged source, so
their complete crossing charts survive, not only their points.

For
distinct partners, choose disjoint source annular collars and a
common identity-type target tube, with exact preimage equal to the
union of those two collars.  Their middle polygons are either both
contractible in the source annulus or both essential.

The mixed alternative is excluded using the original radial map.
A contractible source circle gives a disk filling in the target
stage.  Its partner traverses the same target circle, up to reversal
and a change of initial point.  If that partner were essential in
the source annulus, the strip between it and a marked rim would
give a homotopy to that rim.  Projection and radial retraction would
then contract the once-around loop of $Q^1$, a contradiction.

For each contractible collar let $I_i$ be the closed disk bounded
by its inner polygon and $O_i$ the closed disk bounded by its outer
polygon.  The two $O_i$ are either disjoint, or one lies in the
interior of the other's $I_i$.  This follows from planar separation
and disjointness of the collars; it is a statement about the full
polygon disks, not only their central circles.

If $O_1\subset\operatorname{int}I_0$, choose a finite PL
homeomorphism $H:I_1\to I_0$ with boundary values prescribed by
the common tube parameters.  It is obtained by correcting a disk
homeomorphism on its boundary using the coning extension from
Section~\ref{sec:dehn-proper-disk}.  Define $g=f\circ H^{-1}$ on
$I_0$, retain $f$ outside $\operatorname{int}O_0$, and on the outer
collar use the resolving annulus that joins these two specified
boundary maps.  The entire annular portion between the selected
collars is discarded.  These formulas cover the original source
and agree at the two seams.  The retained source is exactly
$I_1\cup(S\setminus\operatorname{int}O_0)$, where $S$ is the old
annulus.  In the reverse nesting order exchange the two physical
diagonal sheets of the tube and apply the same construction;
the exchange preserves the full longitudinal parameters.

If $O_0$ and $O_1$ are disjoint, choose a boundary-matched finite PL
homeomorphism $H:I_0\to I_1$.  Set $g=f\circ H$ on $I_0$,
$g=f\circ H^{-1}$ on $I_1$, and $g=f$ outside the two open outer
disks.  On the two collars insert the two disjoint resolving annuli.
This exchanges the two inner disk maps once each; it does not make
two copies of either disk.  Their specified boundary values give
the seam identities.  The source remains the same annulus, and
the exact tube preimage gives singleton fibers on the inserted
collars by the same argument as in the reflected case.

In the essential case both collars surround the annulus's hole.
Retain the portion outside the outer selected collar and the portion
inside the inner selected collar but outside the original hole.
Join them with one resolving annulus, discarding the region between
the selected collars.  These two retained pieces and the connecting
strip are annuli joined along complete rims, so their union is an
annulus.  Its boundary maps are homotopic through this union; the
original radial map gives the same orientation on the two prescribed
rims.  The two-rim extension above therefore corrects the whole
parametrization.  In all cases the target map on a retained piece
is literally the old map composed with its source inclusion.

The replacement strips use the same $H_b$ formulas as the interval
surgery, now periodic in their axis coordinate.  The identity closing
identifies each strip with itself, so each replacement is an annulus.
Its full boundary values are used in the finite source attachments.

The count estimate uses these inclusions rather than a drawing of a
smaller singular set.  The finite source-circle decomposition identifies
the intrinsic count with the natural-number cardinality of the connected
components of the actual double-locus set.  The exact tube preimage and
the separated replacement pieces show that every new double component
is carried by an old component outside the selected component or paired orbit.
The selected one is absent.  Rebuilding the local crossing charts
at the joining circles ensures that no component is created there.
This gives a strict inequality of the intrinsic finite component
counts and retains properness, since the tube lies in the interior
and both rims were fixed.  For the reflected case the exact double
relation just proved gives this comparison directly.  For the
paired cases one uses the corresponding retained source inclusions
on the attached finite disks or annuli.  Each old double circle is
either wholly retained or wholly discarded: it cannot cross a
collar boundary, since the only old double points in that collar
lie on its selected middle circle.  The retained set contains neither
selected middle circle.  Partner pairs that lose one member simply
cease to be double pairs.  Thus the new double locus is an embedded
copy of a union of complete old components, with a nonempty selected
orbit removed, which proves the strict finite count inequality.

Strong induction on $c(f)$ now yields an embedded annulus at the
lower stage.  At count zero, any collision would belong to a
nonempty double component; hence the map is injective and compactness
makes it an embedding.  Repeating the proved step backwards along
the finite tower gives an embedded annulus in its initial stage.
The initial open embedding and then the retained inverse chart
preserve its exact rims and whole-frontier properness.

\begin{theorem}[Protected annulus]
\label{thm:dehn-protected-annulus}
\uses{def:dehn-retained-block,thm:dehn-terminal-annulus,thm:dehn-annulus-termination}
For the handle data of Definition~\ref{def:dehn-retained-block} with
$k=1,n=2$, there is a compact embedded chartwise PL annulus
$j:D^1\times Q^1\to R$, with image $T$, such that
\[
 j(\pm1,u)=\pi(\pm1,\tfrac32u),\qquad
 j^{-1}(\partial R)=Q^0\times Q^1,
\]
\[
      T\subset C_2\setminus C_1,
      \qquad T\subset\pi(D^1\times\{y:\|y\|<2\}).
\]
\lean{PoincareMT.M76.Dehn.ProtectedAnnulus.exists_embedded_chart_annulus}
\lean{PoincareMT.M76.Dehn.ProtectedAnnulus.exists_protected_annulus_of_chart}
\source{Hamilton1976}
\dcref{Index-one generalized Dehn application, printed p. 67}
\end{theorem}

\begin{proof}
Use the terminal annulus construction of
Section~\ref{sec:dehn-terminal-annulus} and the descent above.  All
maps remain in the chart domain corresponding to the original strict
shell.  On that shell $e_*^{-1}\circ h=\pi$.  This identity restores
the displayed boundary formulas.  Retained injectivity excludes an
alternative representative in $C_1$, just as in the disk-slab proof.
The strict outer inequality is part of the shell, and compactness
follows from the compact source annulus.
\end{proof}

\section{The same surfaces enclose the protected core}

\begin{theorem}[Termination of marked annulus surgery]
\label{thm:dehn-annulus-termination}
\uses{lem:dehn-annulus-decrease}
\lean{PoincareMT.M76.Dehn.ProtectedAnnulus.exists_embedded_planar_stage_annulus}
A proper finite PL annulus in a retained tower stage, with a finite
ordinary circle decomposition and interior crossing values, admits an
embedded finite PL replacement. The replacement stays in the same
domain, is proper against its whole frontier, and retains both complete
prescribed rim maps.
\end{theorem}

\begin{proof}
Induct on the number of connected components of the actual source
double locus. At zero the map is already embedded. Otherwise
Lemma~\ref{lem:dehn-annulus-decrease} supplies a smaller count and
preserves the hypotheses needed for the next step.
\end{proof}
\label{sec:dehn-enclosing-region}

It remains to find the enclosing region for the surfaces already
constructed.  A new choice of a disk pair or annulus would lose the
markings needed by the later handle argument.

\begin{lemma}[Locally flat sphere separation]
\label{lem:dehn-sphere-separation}
\lean{PoincareMT.M76.LocallyFlatTopologicalSphere.exists_bounded_complement_components}
If
$S\subset E^3$ is the image of a homeomorphism from a two-sphere
and every point has an ambient chart identifying $S$ with a plane,
then $E^3\setminus S=U\sqcup V$ consists of two connected open
sets.  One, $U$, is bounded; its closure is compact and
\[
        \partial U=\partial V=\partial\overline U=S.
\]
\end{lemma}
\begin{proof}
The sphere is two-sided; the local plane charts give a bicollar by
\cite[Theorem 3, p.~339]{Brown1962}. In the one-point compactification,
Brown's complementary-domain theorem gives two domains with common
frontier the sphere \cite[Theorem 5, p.~76]{Brown1960}. Exactly one
contains infinity. The other has compact closure in $\mathbb R^3$,
and the local collar on both sides gives the displayed frontier
equalities. Only separation is used here; the stronger ball recognition
and its relative markings will be used in the rigidity chapter.
\end{proof}

\begin{theorem}[Marked enclosing region]
\label{thm:dehn-enclosing-region}
\uses{thm:dehn-protected-disks,thm:dehn-protected-annulus,lem:dehn-sphere-separation}
Let $T$ be the annulus image of
Theorem~\ref{thm:dehn-protected-annulus}, or the union of the two
disk images of Theorem~\ref{thm:dehn-protected-disks}.
There is a compact region $P\subset R$ such that
\[
 C_1\subset P\subset C_2,\qquad
 P\subset\pi(D^k\times\{y:\|y\|<2\}),\qquad
 C_1\subset\operatorname{int}_R P,
\]
\[
       \partial P=T\cup B_{3/2},\qquad
       P\cap\partial R=B_{3/2}.
\]
The displayed frontier is a chartwise finite PL embedded sphere.
This conclusion specifies an enclosing region; recognition of that
region as a PL ball is a further step in the handle argument.
\lean{PoincareMT.M76.Dehn.ProtectedAnnulus.nonempty_enclosing_region}
\lean{PoincareMT.M76.Dehn.ProtectedDisks.nonempty_enclosing_region}
\source{Hamilton1976}
\dcref{Protected enclosing sphere and compact set, printed p. 67}
\end{theorem}

\begin{proof}
In the annulus case cap the two boundary circles by the attaching
disks $\pi(\{\pm1\}\times\tfrac32D^2)$.  In the disk case join
the disks by $\pi(Q^1\times[-3/2,3/2])$.  Full boundary equality
allows the maps to paste pointwise.  Properness excludes intersections
with the interior of the attaching patch, and simultaneous
disjointness is used in the disk case.  A triangulated prism boundary
therefore parametrizes a finite PL sphere.  At a vertex of a common
subdivision, intersect its local surface star with a small convex
sphere.  The surface link is a simple polygon.  Its two disk sides
on that sphere, coned to the vertex, give an ambient plane chart.
At edge and face points use the corresponding subdivided stars.
This proves local flatness in the retained PL chart.  Pull the
whole sphere and these plane charts back by the original topological
chart inverse; that inverse need not be PL.  The stated separation
result gives its bounded side $U$.  Outside a sufficiently large
coordinate ball the complement is connected, so exactly one of
the two sides is unbounded.  This labels $U$ without making a
choice of a filling ball.

Write $K$ for the source surfaces.  Then
\[
 \partial U=K\cup(Q^{k-1}\times\tfrac32D^n),\quad
 K\subset(D^k\times2D^n)\setminus(D^k\times D^n),\quad
 K\subset D^k\times\{\|y\|<2\}.
\]
Every coordinate of $\overline U$ lies between the extrema of that
coordinate on $\partial U$: a larger value would attain a maximum
in the open set $U$, where moving a little in that coordinate is
possible.  This proves the closed outer bounds.  Since the frontier
is compact and strictly inside the transverse radius-two window,
its transverse coordinate extrema are strictly below two; so are
those of $\overline U$.

The open product $(-1,1)^k\times(-1,1)^n$ is connected and misses
$\partial U$.  It meets $U$, since points of the attaching boundary
patch with transverse coordinate near zero are approached from $U$,
which lies on the domain side of that patch.  It follows that the
whole open product lies in $U$.  Taking closures gives
$D^k\times D^n\subset\overline U$.

At a point of this core on the old boundary, choose a small ball
missing $K$.  Within that ball the frontier of $U$ is contained in
the frontier of the convex domain $D^k\times E^n$.  The intersection
of the ball with the interior of this domain is connected, misses
$\partial U$, and meets $U$.  Hence it lies in $U$.  Density of the
domain interior in the convex closed domain puts the whole relative
ball in $\overline U$.  This proves relative interior containment,
including the old-boundary points of the protected core.  The same
frontier decomposition gives
\[
 \overline U\cap(Q^{k-1}\times E^n)
                 =Q^{k-1}\times\tfrac32D^n.
\]

Finally, the marked quotient is an open embedding on
$E^k\times\{\|y\|<2\}$.  To prove injectivity, any transverse
identification in this window can be paired with first coordinate
zero and tested in the retained closed block; its injectivity
forces the transverse representatives to agree.  Openness comes
from the quotient map.  Set $P=\pi(\overline U)$.  The source region
is compact and lies within this open window.  There is consequently
no extra quotient frontier, and
\[
 \partial\pi(\overline U)=\pi(\partial\overline U).
\]
Restricting the same open projection to the closed domain preserves
relative interiors.  Project all the preceding identities to obtain
the claimed compact region with the original surfaces and markings.
\end{proof}

\section{Protected outputs}
\label{sec:dehn-outputs}

The protected conclusions quantify over the supplied lattice and
atlas before choosing the original chart and its boundary
neighborhood.  In their combined handle application $L$ is a full
discrete lattice; the individual constructions above already work
under the stated discreteness and retained-injectivity conditions.
The generalized Dehn input has two distinct clauses: both protected
surface constructions with their enclosing regions, and the standard
proper-disk theorem for every compact standard PL domain, prescribed
finite PL circle and continuous filling.  The conjunction introduces
no new surface-existence assumption.

For compact-core arguments, the marked loop theorem can be applied
to a relatively open frontier mark and an excluded normal subgroup.
For lower handles, the annulus or disjoint disk pair comes with the
same enclosing sphere and a full attaching-patch equation.  These
are the inputs to protected ball recognition and irreducible
replacement.  The later straightening proof also uses the standard
proper disk in a complementary domain.  Compact-core compression,
prime replacement and relative torus rigidity are separate
constructions; they are not consequences asserted by this chapter.
